\documentclass[10pt,a4paper]{article}
% ============================================================
%  PREAMBLE -- Proteksi Sistem Tenaga Listrik (disalin dari TILI)
% ============================================================
\usepackage[a4paper,margin=2.2cm,top=2.3cm,bottom=2.5cm]{geometry}
\usepackage[T1]{fontenc}
\usepackage[utf8]{inputenc}
\usepackage{lmodern}
\usepackage[indonesian]{babel}
\usepackage{amsmath,amssymb}
\usepackage{icomma}
\usepackage{xcolor}
\usepackage[most]{tcolorbox}
\usepackage[american]{circuitikz}
\usepackage{pgfplots}
\pgfplotsset{compat=1.18}
\usepgfplotslibrary{fillbetween}
\usetikzlibrary{arrows.meta,calc,positioning,decorations.pathreplacing,patterns,angles,quotes,shapes.geometric}
\usepackage{booktabs,array,multirow,longtable,tabularx}
\usepackage{enumitem}
\usepackage{etoolbox}
\usepackage{fontawesome5}
\usepackage{caption}
\captionsetup{font=small,labelfont=bf}
\usepackage{float}
\usepackage[hidelinks]{hyperref}
\setlist{itemsep=1pt,topsep=3pt}
\setlength{\parindent}{0pt}
\setlength{\parskip}{4pt}
\renewcommand{\arraystretch}{1.18}

% ------------------------------------------------------------
%  PALET: 5 keluarga warna x 4 gradasi.
%  Makna tiap warna DITETAPKAN ULANG di awal setiap bab.
% ------------------------------------------------------------
\definecolor{biru1}{HTML}{0D47A1}  % biru tua
\definecolor{biru2}{HTML}{1E88E5}  % biru terang
\definecolor{biru3}{HTML}{00838F}  % biru kehijauan (sian)
\definecolor{biru4}{HTML}{3949AB}  % biru nila
\definecolor{merah1}{HTML}{B71C1C} % merah tua
\definecolor{merah2}{HTML}{E53935} % merah terang
\definecolor{merah3}{HTML}{EF6C00} % jingga
\definecolor{merah4}{HTML}{AD1457} % merah keunguan
\definecolor{hijau1}{HTML}{1B5E20} % hijau tua
\definecolor{hijau2}{HTML}{43A047} % hijau terang
\definecolor{hijau3}{HTML}{827717} % hijau zaitun
\definecolor{hijau4}{HTML}{00796B} % hijau toska
\definecolor{ungu1}{HTML}{4A148C}  % ungu tua
\definecolor{ungu2}{HTML}{8E24AA}  % ungu terang
\definecolor{ungu3}{HTML}{5E35B1}  % ungu kebiruan
\definecolor{ungu4}{HTML}{C2185B}  % ungu kemerahan (magenta)
\definecolor{coklat1}{HTML}{5D4037} % coklat tua
\definecolor{coklat2}{HTML}{A1572A} % coklat bata
\definecolor{coklat3}{HTML}{B8860B} % kuning tua / emas
\definecolor{coklat4}{HTML}{8D6E63} % coklat muda

% ------------------------------------------------------------
%  Simbol berwarna
%  \defq{nama}{warna}{simbol}{nama besaran}
%    -> \nama   : simbol berwarna (bisa di teks maupun rumus)
%    -> \tnama  : nama besaran, tebal & berwarna sama
% ------------------------------------------------------------
\newcommand{\defsym}[3]{\expandafter\gdef\csname #1\endcsname{\ensuremath{\textcolor{#2}{#3}}}}
\newcommand{\kw}[2]{\textcolor{#1}{\textbf{#2}}}
\newcommand{\defq}[4]{\defsym{#1}{#2}{#3}\expandafter\gdef\csname t#1\endcsname{\kw{#2}{#4}}}
\newcommand{\cl}[2]{\textcolor{#1}{#2}}   % warnai potongan rumus apa saja

% ------------------------------------------------------------
%  Kotak-kotak
% ------------------------------------------------------------
% Rumus kunci: bingkai tipis
\newtcbox{\kotak}{enhanced,colback=white,colframe=black!60,boxrule=0.5pt,arc=0pt,
  left=6pt,right=6pt,top=3pt,bottom=3pt,nobeforeafter,tcbox raise base}
% Setiap kotak rumus bernomor (bab.urutan), misalnya (2.5). Pasang \label{r:...} tepat sesudahnya.
\newcounter{rumus}[section]
\renewcommand{\therumus}{\thesection.\arabic{rumus}}
\newcommand{\nomorwarna}{biru1}
\newsavebox{\rumusbox}
\newcommand{\rumus}[2][]{\par\refstepcounter{rumus}\sbox{\rumusbox}{\kotak{%
  {\sffamily\bfseries\color{\nomorwarna}(\therumus)}\quad
  \if\relax\detokenize{#1}\relax\else{\scriptsize\sffamily\color{black!55}#1}\quad\fi
  $\displaystyle #2$}}\begin{center}\ifdim\wd\rumusbox>\linewidth
  \resizebox{\linewidth}{!}{\usebox{\rumusbox}}\else\usebox{\rumusbox}\fi\end{center}\par}

% Pemahaman penting: warna mengikuti keluarga konsepnya
\newtcolorbox{paham}[2][biru1]{enhanced,breakable,colback=#1!5,colframe=#1,
  boxrule=0pt,leftrule=3pt,arc=0pt,outer arc=0pt,
  colbacktitle=#1!13,coltitle=#1!80!black,fonttitle=\bfseries,
  title={\faLightbulb[regular]\enspace Pemahaman penting: #2},
  left=7pt,right=7pt,top=4pt,bottom=4pt,before skip=8pt,after skip=8pt}

% Proses di balik layar (penurunan / alasan): netral abu-abu
\newtcolorbox{balik}[1]{enhanced,breakable,colback=black!2.5,colframe=black!45,
  boxrule=0pt,leftrule=3pt,arc=0pt,outer arc=0pt,
  colbacktitle=black!8,coltitle=black!80,fonttitle=\bfseries,
  title={\faCogs\enspace Proses di balik layar: #1},
  left=7pt,right=7pt,top=4pt,bottom=4pt,before skip=8pt,after skip=8pt}

% Koreksi / catatan tentang slide: garis putus-putus
\newtcolorbox{koreksi}[1][Catatan tentang slide]{enhanced,breakable,colback=white,
  frame hidden,borderline={0.7pt}{0pt}{black!65,dashed},arc=0pt,
  fonttitle=\bfseries,coltitle=black,colbacktitle=white,
  title={\faExclamationTriangle\enspace #1},
  left=7pt,right=7pt,top=3pt,bottom=4pt,before skip=8pt,after skip=8pt}

% Contoh soal
\newtcolorbox{contoh}[1]{enhanced,breakable,colback=white,colframe=black!40,
  boxrule=0.5pt,arc=1pt,colbacktitle=black!8,coltitle=black,fonttitle=\bfseries,
  title={\faEdit[regular]\enspace Contoh Soal: #1},
  left=7pt,right=7pt,top=4pt,bottom=4pt,before skip=8pt,after skip=8pt}
\newcommand{\jawab}{\par\smallskip\textbf{Penyelesaian.}\ }

% Prosedur penyelesaian: Diketahui -> Ditanya -> Rencana -> Langkah 1, 2, ...
% Setiap langkah: \cari{besaran} lalu alasan (kenapa sekarang), rumus dalam simbol, baru angka.
\newcommand{\prosedurwarna}{biru3!75!black}
% Diketahui: butir dipisah titik koma (;), tiap butir satu baris.
\DeclareListParser*{\dkparse}{;}
\newcommand{\dkitem}[1]{\par\hangindent=2.2em\noindent\makebox[2.2em][l]{\hspace{0.9em}\textcolor{\prosedurwarna}{\textbullet}}\ignorespaces#1}
\newcommand{\diket}[1]{\par\smallskip\noindent\textbf{\textcolor{\prosedurwarna}{Diketahui.}}\par
  {\raggedright\dkparse{\dkitem}{#1}\par}}
\newcommand{\ditanya}[1]{\par\noindent\textbf{\textcolor{\prosedurwarna}{Ditanya.}}\ #1\par}
\newcommand{\rencana}[1]{\par\noindent\textbf{\textcolor{\prosedurwarna}{Rencana.}}\ #1\par}
\newcommand{\bagsoal}[1]{\par\smallskip\noindent\textbf{#1}\par}
% Langkah: judul "Langkah n. Cari X." lalu baris-baris berlabel di bawahnya.
\newcounter{lkh}
\newenvironment{langkah}{\par\setcounter{lkh}{0}\raggedright
  \setlength{\lineskiplimit}{3pt}\setlength{\lineskip}{4pt}
  \def\item{\par\addvspace{5pt}\stepcounter{lkh}\noindent
    {\bfseries\color{\prosedurwarna}Langkah \arabic{lkh}.}\ }}{\par\smallskip}
\newcommand{\cari}[1]{\textbf{Cari #1.}\ }
\newcommand{\cek}{\textbf{Cek.}\ }
\newlength{\lblw}\setlength{\lblw}{5.4em}
\newcommand{\baris}[2]{\par\hangindent=\lblw\hangafter=1\noindent
  \makebox[\lblw][l]{\hspace{1.3em}{\small\sffamily\color{black!55}#1}}#2\par}
\newcommand{\alasan}[1]{\baris{Kenapa}{#1}}
% rumus dari kotak bernomor: nomor + rumus ditulis ulang
\newcommand{\pakai}[2]{\baris{Rumus}{{\sffamily\bfseries\color{\nomorwarna}(\ref{#1})}\enspace$\displaystyle #2$}}
% rumus dasar yang tidak berasal dari kotak
\newcommand{\pakaidasar}[2]{\baris{Rumus}{{\small\sffamily\color{black!60}[#1]}\enspace$\displaystyle #2$}}
\newcommand{\susun}[2][]{\baris{Susun}{$\displaystyle #2$\if\relax\detokenize{#1}\relax\else
  \enspace{\small\color{black!60}(#1)}\fi}}
\newcommand{\angka}[1]{\baris{Angka}{$\displaystyle #1$}}
\newcommand{\hasil}[1]{\baris{Hasil}{\colorbox{biru3!10}{$\displaystyle #1$}}}
\newcommand{\hasilt}[1]{\baris{Hasil}{#1}}
\newcommand{\ceks}[1]{\baris{Cek}{#1}}

% Peta warna bab
\newcommand{\swatch}[1]{\textcolor{#1}{\rule{0.9em}{0.9em}}}
\newcommand{\swatches}[1]{\foreach \c in {#1}{\swatch{\c}\,}}
\newtcolorbox{petawarna}[1]{enhanced,colback=white,colframe=black!25,boxrule=0.4pt,arc=1pt,
  title={\faPalette\enspace Peta warna #1 \ (berlaku hanya di bab ini)},
  colbacktitle=black!6,coltitle=black,fonttitle=\bfseries\small,
  left=4pt,right=4pt,top=2pt,bottom=2pt,fontupper=\small}
\newenvironment{pw}{\begin{tabular}{@{}l>{\raggedright\arraybackslash}p{0.32\linewidth}>{\raggedright\arraybackslash}p{0.47\linewidth}@{}}}{\end{tabular}}
\newcommand{\pwrow}[3]{\swatches{#1} & \textbf{#2} & #3 \\}

% Kamus variabel
\newenvironment{kamus}{\par\medskip\textbf{Kamus variabel}\par\nopagebreak\small
  \setlength{\LTpre}{3pt}\setlength{\LTpost}{6pt}
  \begin{longtable}{@{}>{\centering\arraybackslash}p{2.0cm}>{\raggedright\arraybackslash}p{\dimexpr\linewidth-4.8cm-4\tabcolsep\relax}>{\raggedright\arraybackslash}p{2.8cm}@{}}
  \toprule Simbol & Nama dan arti & Satuan\\ \midrule \endhead
  \bottomrule \endfoot}{\end{longtable}\normalsize}

% Bab baru selalu di halaman baru
\newcommand{\bab}[1]{\clearpage\section{#1}}

% Panah fasor
\tikzset{fasor/.style={-{Stealth[length=2.6mm]},line width=1.1pt},
         bantu/.style={densely dashed,black!45}}

% Tata letak seksi
\usepackage{titlesec}
\titleformat{\section}{\Large\bfseries}{\thesection}{0.8em}{}
\titleformat{\subsection}{\large\bfseries}{\thesubsection}{0.7em}{}
\titleformat{\subsubsection}{\normalsize\bfseries}{\thesubsubsection}{0.6em}{}
\titlespacing*{\section}{0pt}{4pt}{10pt}

% ------------------------------------------------------------
%  Tambahan (dari TILI, dipakai juga di Proteksi)
% ------------------------------------------------------------
\usepackage{graphicx}
\graphicspath{{img/}}
\usepackage{needspace}
\usetikzlibrary{fit,backgrounds,shapes.misc,shapes.symbols}
\newcommand{\fig}[1]{\input{fig/#1}}
% beberapa foto sejajar: \fotos{tinggi}{file1/ket1,file2/ket2,...}
\newcommand{\fotos}[2]{\par\begin{center}\foreach \f/\k in {#2}{\begin{tabular}[b]{@{}c@{}}\includegraphics[height=#1]{\f}\\[-1pt]
  {\footnotesize\color{black!60}\emph{\k}}\end{tabular}\hspace{6mm}}\end{center}\par}
% gambar di tengah, diperkecil otomatis bila lebih lebar dari baris
\newsavebox{\figbox}
\newcommand{\figc}[1]{\par\begin{center}\sbox{\figbox}{\input{fig/#1}}\ifdim\wd\figbox>\linewidth
  \resizebox{\linewidth}{!}{\usebox{\figbox}}\else\usebox{\figbox}\fi\end{center}\par}
% foto slide: \foto[lebar]{file}{keterangan}
\newcommand{\foto}[3][0.7\linewidth]{\par\begin{center}\begin{minipage}{\linewidth}\centering
  \includegraphics[width=#1,height=5.6cm,keepaspectratio]{#2}\par\vspace{1pt}
  {\footnotesize\color{black!60}\emph{#3}}\end{minipage}\end{center}\par}
% gaya node diagram
\newcommand{\nohyph}{\hyphenpenalty=10000\exhyphenpenalty=10000\relax}
\tikzset{
  blok/.style={draw=black!60,rounded corners=2pt,align=center,font=\footnotesize,minimum height=7mm,inner sep=3pt,fill=black!3,execute at begin node=\nohyph},
  ar/.style={-{Stealth[length=2.2mm]},black!65,thick},
  bus/.style={line width=2.2pt},
  lbl/.style={font=\scriptsize,align=center,execute at begin node=\nohyph},
}
% simbol kecil untuk diagram satu garis
\newcommand{\cbsym}[2]{\draw[thick,fill=white] (#1) ++(-0.13,-0.13) rectangle ++(0.26,0.26); \node[font=\tiny,right] at ($(#1)+(0.15,0)$) {#2};}
\newcommand{\trafosym}[1]{\draw[thick] ($(#1)+(0,0.17)$) circle (0.2); \draw[thick] ($(#1)+(0,-0.17)$) circle (0.2);}

% ------------------------------------------------------------
%  Tambahan khusus Proteksi
% ------------------------------------------------------------
\usepgfplotslibrary{groupplots}
% simbol sekering kecil (zig-zag) di garis horizontal: \fusesym{titik}
\newcommand{\fusesym}[1]{\draw[thick] ($(#1)+(-0.22,0)$) .. controls ($(#1)+(-0.11,0.32)$) and ($(#1)+(-0.02,0.32)$) .. ($(#1)$) .. controls ($(#1)+(0.02,-0.32)$) and ($(#1)+(0.11,-0.32)$) .. ($(#1)+(0.22,0)$);}
% sekering pada garis vertikal
\newcommand{\fusesymv}[1]{\draw[thick] ($(#1)+(0,0.22)$) .. controls ($(#1)+(0.32,0.11)$) and ($(#1)+(0.32,0.02)$) .. ($(#1)$) .. controls ($(#1)+(-0.32,-0.02)$) and ($(#1)+(-0.32,-0.11)$) .. ($(#1)+(0,-0.22)$);}
% gelembung arus gangguan max/min
\newcommand{\gel}[3]{\node[draw,ellipse,inner sep=0.6pt,font=\tiny,align=center] at (#1) {#2\\[-1pt]\rule{4.5mm}{0.3pt}\\[-1pt]#3};}
% Pelengkap di luar slide: latar abu-abu, garis kiri tipis
\newtcolorbox{luar}{enhanced,breakable,colback=black!3,colframe=coklat4,
  boxrule=0pt,leftrule=1.5pt,arc=0pt,outer arc=0pt,fontupper=\small,
  title={\faBookOpen\enspace Pelengkap (di luar slide)},colbacktitle=black!3,coltitle=coklat1,
  fonttitle=\bfseries\small,left=7pt,right=7pt,top=3pt,bottom=3pt,before skip=6pt,after skip=6pt}
% kilat gangguan
\newcommand{\kilat}[1]{\draw[merah1,thick,-{Stealth[length=1.6mm]}] ($(#1)+(0.08,0.30)$)--($(#1)+(-0.06,0.05)$)--($(#1)+(0.08,0.02)$)--($(#1)+(-0.05,-0.26)$);}

\newgeometry{a4paper,margin=1.7cm,top=1.6cm,bottom=1.9cm}
\setlength{\parskip}{3pt}
\hypersetup{pdftitle={Ringkasan Metode GEMASTIK NEW-GEN}}
\usepackage{listings}
\newcommand{\rfvalsatu}{79{,}4}
\newcommand{\rfvalbanyak}{87{,}4}
\newcommand{\lrwa}{3{,}65}
\newcommand{\lrwb}{3{,}74}
\newcommand{\lrc}{0{,}74}
\newcommand{\lracc}{93{,}3}
\newcommand{\kmiter}{4}
\newcommand{\pcasatu}{96}
\newcommand{\pcadua}{4}
\newcommand{\ifluar}{2{,}9}
\newcommand{\iftengah}{11{,}5}
\newcommand{\esmaekecil}{9{,}2}
\newcommand{\esmaebesar}{8{,}5}
\newcommand{\mfmu}{3{,}19}
\newcommand{\mfbu}{-0{,}28}
\newcommand{\mfbi}{-0{,}15}
\newcommand{\mfpq}{-1{,}36}
\newcommand{\mfhasil}{1{,}39}
\newcommand{\qgrup}{0{,}551}
\newcommand{\qsatu}{0{,}000}
\newcommand{\jumlahsisi}{35}


% ---------- kotak versi ringkas (disalin dari Ringkasan Proteksi) ----------
\newtcolorbox{inti}[1][biru1]{enhanced,breakable,colback=#1!6,colframe=#1,boxrule=0pt,leftrule=3pt,arc=0pt,
  outer arc=0pt,left=6pt,right=6pt,top=2pt,bottom=2pt,before skip=5pt,after skip=5pt,
  fontupper=\small,before upper={\textcolor{#1}{\faLightbulb[regular]\ \textbf{Inti.}}\ }}
\newtcolorbox{kenapa}[1][]{enhanced,breakable,colback=black!3,colframe=black!40,boxrule=0pt,leftrule=3pt,arc=0pt,
  outer arc=0pt,left=6pt,right=6pt,top=2pt,bottom=2pt,before skip=5pt,after skip=5pt,
  fontupper=\small,before upper={\textcolor{black!70}{\faCogs\ \textbf{Kenapa begitu?}\if\relax\detokenize{#1}\relax\else\ (#1)\fi}\ }}
\newtcolorbox{jebak}{enhanced,breakable,colback=white,frame hidden,borderline={0.6pt}{0pt}{merah1!70,dashed},arc=0pt,
  left=6pt,right=6pt,top=2pt,bottom=2pt,before skip=5pt,after skip=5pt,fontupper=\small,
  before upper={\textcolor{merah1}{\faExclamationTriangle\ \textbf{Jebakan lomba.}}\ }}
\newtcolorbox{soal}[1]{enhanced,breakable,colback=white,colframe=black!35,boxrule=0.4pt,arc=1pt,
  colbacktitle=black!7,coltitle=black,fonttitle=\bfseries\small,title={\faCode\ Di notebook: #1},
  left=6pt,right=6pt,top=2pt,bottom=2pt,before skip=5pt,after skip=5pt,fontupper=\small}
\newtcolorbox[auto counter]{diag}[2][]{enhanced,breakable,colback=white,colframe=biru3!70!black,boxrule=0.5pt,arc=1pt,
  colbacktitle=biru3!12,coltitle=black,fonttitle=\bfseries\small,title={\faPen\ Diagram \thetcbcounter: #2},label={#1},
  left=6pt,right=6pt,top=3pt,bottom=3pt,before skip=6pt,after skip=6pt,fontupper=\small}
\renewtcolorbox{luar}{enhanced,breakable,colback=white,frame hidden,borderline={0.6pt}{0pt}{black!55,dotted},arc=0pt,
  left=6pt,right=6pt,top=2pt,bottom=2pt,before skip=5pt,after skip=5pt,fontupper=\small,
  before upper={\textcolor{black!70}{\faBook\ \textbf{Kerabatnya.}}\ }}
\renewcommand{\rumus}[2][]{\par\vspace{-2pt}\begin{center}\kotak{%
  \if\relax\detokenize{#1}\relax\else{\scriptsize\sffamily\color{black!55}#1}\quad\fi
  $\displaystyle #2$}\end{center}\vspace{-4pt}\par}
\newcolumntype{P}[1]{>{\raggedright\arraybackslash}p{#1}}

\newcommand{\rb}[3]{\needspace{12\baselineskip}\section{#1}\vspace{-6pt}{\footnotesize\textcolor{black!55}{\faPalette\ Warna bab ini:}\ #2}\par
  \vspace{1pt}{\small\emph{#3}}\par}
\newcommand{\lgd}[2]{\swatch{#1}\,#2\quad}
\newcommand{\sub}[1]{\par\smallskip{\normalsize\bfseries #1}\par\nopagebreak}
\titlespacing*{\section}{0pt}{8pt}{4pt}
\setlist{itemsep=0pt,topsep=2pt,leftmargin=*}
\newcommand{\figs}[2][0.8]{\par\begin{center}\sbox{\figbox}{\input{fig/#2}}\ifdim\wd\figbox>#1\linewidth
  \resizebox{#1\linewidth}{!}{\usebox{\figbox}}\else\usebox{\figbox}\fi\end{center}\par}

% kode Python bergaya notebook
\lstset{language=Python,basicstyle=\ttfamily\scriptsize,commentstyle=\color{hijau1},keywordstyle=\color{biru1},
  stringstyle=\color{merah1},columns=fullflexible,keepspaces=true,showstringspaces=false,breaklines=true,
  aboveskip=2pt,belowskip=2pt,upquote=true}
\newcommand{\kode}[1]{\texttt{\small #1}}

\begin{document}

% ============================================================ JUDUL + CARA PAKAI
\begin{center}
{\LARGE\bfseries Ringkasan Metode GEMASTIK NEW-GEN}\\[3pt]
{\normalsize Model pohon, linear, jarak, neural network, tanpa label, teks, deret waktu, rekomendasi, graf, dan cara menilai model}\\[4pt]
\begin{tikzpicture}
  \foreach \c [count=\i] in {biru1,biru2,biru3,biru4,merah1,merah2,merah3,merah4,hijau1,hijau2,hijau3,hijau4,ungu1,ungu2,ungu3,ungu4,coklat1,coklat2,coklat3,coklat4}
    \fill[\c] ({0.42*\i},0) rectangle ++(0.38,0.22);
\end{tikzpicture}
\end{center}

\begin{tcolorbox}[enhanced,colback=black!2,colframe=black!30,boxrule=0.4pt,arc=1pt,fontupper=\small,
  title={Cara memakai ringkasan ini},colbacktitle=black!7,coltitle=black,fonttitle=\bfseries\small]
\begin{itemize}
\item Sumber: notebook di \kode{\_.NEW-GEN/} (199 notebook) dan skor di \kode{KATALOG.md}. Semua gambar dibuat dari
\textbf{data mainan} dengan \kode{buat\_gambar.py}, jadi angkanya nyata tetapi bukan data lomba.
\item \textbf{Warna direset tiap bab.} Baris kecil di bawah judul bab menjelaskan arti warnanya. Di gambar 2D:
\textcolor{biru2}{\textbf{biru}} = kelas 0, \textcolor{merah3}{\textbf{jingga}} = kelas 1, warna latar = tebakan model
(makin pekat = makin yakin).
\item Kotak: \textcolor{biru1}{\textbf{Inti}} = ide yang membuat hal lain masuk akal; \kotak{rumus} = cukup paham artinya;
\textbf{Kenapa begitu?} = alasan bertahap; \textcolor{biru3!70!black}{\textbf{Diagram}} = gambar + cara membaca;
\textcolor{merah1}{\textbf{Jebakan lomba}} = kesalahan yang sering terjadi; \textbf{Di notebook} = baris kode yang
benar-benar ada di toolkit; \textbf{Kerabatnya} = metode lain di keluarga yang sama.
\item \textbf{Aturan emas:} model dinilai pada data yang \emph{belum pernah dilihat}. Semua teknik di sini (split, KFold,
\kode{max\_depth}, \kode{learning\_rate}, regularisasi) adalah cara menjamin skor validasi jujur dan model tidak menghafal.
\end{itemize}
\end{tcolorbox}
{\small\tableofcontents}

% ============================================================ BAB 1
\rb{Peta Besar: Dari Soal ke Metode}
{\lgd{biru1}{menyiapkan data}\lgd{hijau1}{model}\lgd{merah1}{menilai}}
{Setiap notebook NEW-GEN punya bentuk yang sama: 8 sel berurutan. Yang berubah antar-soal hanya isi sel Model.}

\begin{diag}[dg:alur]{Alur satu notebook}
\figs[0.95]{alur}
Sel 1--4 \textcolor{biru1}{menyiapkan data}, sel 5--6 \textcolor{hijau1}{melatih model}, sel 7--8 \textcolor{merah1}{menilai
dan mengirim}. \textbf{Cara membaca:} panah putus-putus = putaran percobaan; ganti model, jalankan ulang, bandingkan skor.
\end{diag}

\begin{inti}[hijau1]
Satu notebook = satu metode. Di sel Model ada \textbf{satu baris} \kode{model = ...} yang aktif, dan alternatifnya ditulis di
baris \kode{\#}. Mencoba model lain = pindahkan tanda \kode{\#}, lalu \emph{Run All}. Skor tiap alternatif pada data demo
ada di \kode{KATALOG.md}.
\end{inti}

\sub{Soal jenis apa, coba metode apa}
{\small
\begin{tabularx}{\linewidth}{@{}P{3.0cm}P{3.0cm}XP{3.6cm}@{}}
\toprule
Jenis soal & Contoh & Model aktif $\to$ alternatif kuat & Notebook\\ \midrule
Tabel $\to$ kelas & pelanggan berhenti? & LogisticRegression $\to$ LightGBM, XGBoost, CatBoost & \kode{tabular/klasifikasi}\\
Tabel $\to$ angka & harga rumah & LightGBM $\to$ XGBoost, Ridge & \kode{tabular/regresi}\\
Teks $\to$ kelas & sentimen ulasan & TF-IDF + LinearSVC $\to$ embedding, fine-tune IndoBERT & \kode{text/klasifikasi\_tfidf}\\
Gambar $\to$ kelas & citra medis & CNN pretrained (ResNet) & \kode{medis/klasifikasi\_citra}\\
Cari objek di gambar & hitung kendaraan & YOLO & \kode{vision/deteksi}\\
Ramal masa depan & penjualan bulan depan & fitur lag + RandomForest $\to$ LightGBM; SARIMAX & \kode{timeseries/peramalan\_ml}\\
Kelompok tanpa label & segmen pelanggan & KMeans (+ PCA untuk gambar) & \kode{unsupervised/clustering}\\
Cari yang aneh & transaksi curang & IsolationForest $\to$ LOF, ECOD & \kode{timeseries/anomali}\\
Rekomendasi & tebak rating film & SVD (matrix factorization) & \kode{recsys/prediksi\_rating}\\
Jaringan / graf & komunitas & Louvain & \kode{graph/komunitas}\\ \bottomrule
\end{tabularx}}

\sub{Empat konsep dasar}
{\small
\begin{tabularx}{\linewidth}{@{}>{\bfseries\raggedright\arraybackslash}p{3.2cm}X@{}}
\toprule
\cl{biru1}{Fitur dan label} & $X$ = kolom yang kita tahu (umur, gaji, jumlah keluhan); $y$ = yang ingin ditebak (berhenti
atau tidak). Model mencari hubungan $X \to y$.\\
\cl{biru1}{Klasifikasi vs regresi} & klasifikasi menebak \textbf{kelas}; regresi menebak \textbf{angka}. Banyak model punya dua
versi: \kode{XGBClassifier} dan \kode{XGBRegressor}.\\
\cl{merah1}{Latih vs validasi} & sebagian data disembunyikan dari model; skor dihitung di sana karena di lomba model juga dinilai
pada data yang belum pernah dilihat.\\
\cl{merah1}{Overfitting} & model terlalu rumit $\to$ menghafal data latih termasuk noise-nya: skor latih tinggi, skor validasi
turun (Diagram~\ref{dg:treekurva}).\\ \bottomrule
\end{tabularx}}

\begin{jebak}
\textbf{Kebocoran data (leakage).} Skor validasi tinggi sekali tetapi skor lomba jeblok, biasanya karena informasi dari data
validasi/test bocor ke latihan: \kode{StandardScaler} di-\kode{fit} pada seluruh data, kolom yang sebenarnya hasil dari target
(misal ``tanggal\_berhenti'' untuk menebak berhenti), atau deret waktu yang diacak. Karena itu toolkit memakai \kode{Pipeline}:
preprocessing ikut di-\kode{fit} hanya pada data latih.
\end{jebak}

% ============================================================ BAB 2
\rb{Model Berbasis Pohon}
{\lgd{biru2}{kelas 0}\lgd{merah3}{kelas 1}\lgd{biru1}{tebakan model / validasi}\lgd{merah1}{sisa kesalahan / latih}}
{Keluarga paling sering menang di data tabel. Semuanya dibangun dari satu bahan: pohon keputusan.}

\subsection{Decision Tree}
\begin{inti}[biru1]
Model bertanya ya/tidak berulang kali, misalnya ``umur $<$ 30?'' lalu ``gaji $<$ 5 juta?'', sampai isi tiap kotak cukup
seragam. Setiap pertanyaan memotong ruang data dengan \textbf{garis lurus tegak atau mendatar}. Tebakan = kelas mayoritas
di kotak tempat titik jatuh.
\end{inti}

\begin{diag}[dg:tree]{Pohon dangkal vs pohon dalam}
\figs[1.0]{tree}
Garis tipis = potongan pohon. \textbf{Cara membaca:} kiri, paling banyak 3 pertanyaan berurutan, kotaknya besar dan agak kasar.
Kanan, sampai 12 pertanyaan berurutan, banyak kotak kecil yang dibuat hanya untuk satu-dua titik (yang sebenarnya noise).
Akurasi latih naik ke 100\,\%, tetapi akurasi validasi justru turun.
\end{diag}

\rumus[Gini (ukuran ``campur aduk'')]{\text{Gini} = 1 - \sum_k p_k^2 \qquad (0 = \text{murni satu kelas},\ 0{,}5 = \text{campur rata 2 kelas})}

\begin{kenapa}[cara memilih potongan]
(1) Di tiap kotak, model mencoba \emph{semua} garis yang mungkin pada \emph{semua} fitur. (2) Untuk tiap garis, hitung Gini
kiri dan kanan (dibobot jumlah titik). (3) Pilih garis yang Gini-nya paling kecil. (4) Ulangi di tiap kotak baru sampai
\kode{max\_depth} tercapai atau kotak sudah murni. Tidak ada rumus ``garis terbaik''; pohon mencoba satu per satu.
\end{kenapa}

\begin{diag}[dg:treekurva]{Overfitting terlihat di kurva}
\figs[0.6]{treekurva}
\textbf{Cara membaca:} garis \textcolor{merah1}{merah} (latih) terus naik mendekati 100\,\% karena pohon makin pandai menghafal.
Garis \textcolor{biru1}{biru} (validasi) berhenti naik lalu turun. Kedalaman terbaik = puncak garis biru, bukan garis merah.
\end{diag}

\subsection{Random Forest dan ExtraTrees}
\begin{inti}[biru1]
Buat ratusan pohon yang dalam. Tiap pohon hanya melihat \textbf{sampel acak} dari baris data (bootstrap) dan \textbf{fitur acak}
di tiap potongan, jadi kesalahannya berbeda-beda. Tebakan akhir = \textbf{rata-rata suara} semua pohon.
\end{inti}

\begin{diag}[dg:forest]{Satu pohon vs hutan 50 pohon}
\figs[1.0]{forest}
\textbf{Cara membaca:} satu pohon dalam (kiri) membuat batas kotak-kotak tajam dan percaya diri walaupun salah. Rata-rata 50 pohon
(kanan) membuat batas halus; warna pudar di perbatasan = pohon-pohon tidak sepakat. Validasi naik dari \rfvalsatu\,\% ke
\rfvalbanyak\,\%.
\end{diag}

\begin{kenapa}[kenapa rata-rata pohon lebih baik]
Satu pohon dalam punya \emph{variansi} tinggi: ganti sedikit datanya, bentuknya berubah jauh. Kalau kesalahan tiap pohon
acak dan tidak searah, merata-ratakan banyak pohon membuat kesalahan saling meniadakan, sedangkan pola yang benar (yang
dilihat hampir semua pohon) tetap tinggal. Mirip bertanya ke 500 orang yang masing-masing membaca sebagian berita.
\end{kenapa}

\begin{luar}
\kode{ExtraTrees} = Random Forest yang titik potongnya juga diacak (tidak dicari yang terbaik), jadi lebih cepat dan kadang lebih
tahan noise. Kenop: \kode{n\_estimators} (300--500 cukup; lebih banyak tidak merusak, hanya lebih lambat), \kode{max\_features}.
\end{luar}

\subsection{Gradient Boosting: XGBoost, LightGBM, CatBoost}
\begin{inti}[merah1]
Pohon dibuat \textbf{satu per satu}. Pohon pertama menebak kasar. Pohon kedua \emph{tidak} menebak target, tetapi menebak
\textbf{sisa kesalahan} (residual) pohon pertama. Pohon ketiga menebak sisa kesalahan berikutnya, dan seterusnya. Tebakan akhir
= jumlah semua pohon. Setiap pohon dikali \kode{learning\_rate} supaya langkahnya kecil dan hati-hati.
\end{inti}

\rumus[boosting]{F_m(x) = F_{m-1}(x) + \eta\, h_m(x), \qquad h_m \text{ dilatih menebak } y - F_{m-1}(x)}

\begin{diag}[dg:boost]{Boosting ronde demi ronde (regresi 1 fitur)}
\figs[1.0]{boost}
\textbf{Baris atas:} titik abu = data; garis \textcolor{biru1}{biru} = jumlah semua pohon sejauh ini; garis tegak
\textcolor{merah1}{merah} = sisa kesalahan tiap titik. \textbf{Baris bawah:} sisa kesalahan itu digambar sendiri (titik merah),
dan garis tangga \textcolor{merah3}{jingga} = pohon kecil berikutnya yang dilatih menebaknya. Dari 1 ke 5 ke 50 pohon, garis
biru makin menempel ke pola data dan sisa kesalahan makin acak (tidak berpola lagi).
\end{diag}

\begin{diag}[dg:boostkurva]{Learning rate besar vs kecil}
\figs[0.62]{boostkurva}
\textbf{Cara membaca:} putus-putus = latih, garis penuh = validasi. \kode{learning\_rate} 0,3 (merah) turun cepat lalu validasinya
mulai naik lagi (overfitting). 0,05 (biru) turun pelan, butuh lebih banyak pohon, tetapi titik terendah validasinya
setara atau lebih baik dan lebih stabil.
\end{diag}

\begin{kenapa}[kenapa menebak sisa kesalahan]
(1) Pohon dangkal sendirian lemah, hanya bisa membuat 4--8 anak tangga. (2) Kalau pohon berikutnya fokus ke bagian yang masih
salah, kemampuan bertambah tepat di tempat yang dibutuhkan. (3) Menjumlahkan ratusan tangga kecil menghasilkan bentuk apa pun.
(4) \kode{learning\_rate} kecil = tiap pohon hanya memperbaiki sebagian, sehingga satu pohon yang keliru tidak merusak semuanya.
Nama ``gradient'' karena sisa kesalahan $y - F$ adalah arah turun (gradien) dari error kuadrat.
\end{kenapa}

{\small
\begin{tabularx}{\linewidth}{@{}P{2.6cm}XXXX@{}}
\toprule
 & \textbf{XGBoost} & \textbf{LightGBM} & \textbf{CatBoost} & \textbf{HistGradient\-Boosting}\\ \midrule
Cara tumbuh & lapis demi lapis, rata & daun demi daun: yang paling untung dipecah dulu & pohon simetris & mirip LightGBM\\
Kelebihan & stabil, paling populer & sangat cepat di data besar & kolom kategori tanpa encoding & bawaan sklearn\\
Kenop & \kode{n\_estimators}, \kode{learning\_rate}, \kode{max\_depth} & \kode{n\_estimators}, \kode{learning\_rate},
\kode{num\_leaves} & \kode{iterations}, \kode{learning\_rate}, \kode{depth} & \kode{max\_iter}, \kode{learning\_rate}\\ \bottomrule
\end{tabularx}}

\begin{soal}{tabular/regresi (baris model aktif)}
\begin{lstlisting}
model = LGBMRegressor(n_estimators=1000, learning_rate=0.03, random_state=SEED, verbose=-1)
# model = XGBRegressor(n_estimators=500, learning_rate=0.03, random_state=SEED)
# model = CatBoostRegressor(iterations=500, random_seed=SEED, verbose=0)
\end{lstlisting}
Aturan praktis: \kode{learning\_rate} 0,03--0,1; \kode{n\_estimators} 300--1000; pohon dangkal (depth 3--8 atau
\kode{num\_leaves} 15--63). Turunkan \kode{learning\_rate} $\Rightarrow$ naikkan \kode{n\_estimators}.
\end{soal}

\begin{jebak}
Random Forest dan boosting \emph{terdengar} mirip karena sama-sama ``banyak pohon''. Bedanya: Random Forest = pohon
\textbf{bersamaan, independen, dirata-rata} (menambah pohon tidak bikin overfitting). Boosting = pohon \textbf{berurutan, saling
memperbaiki, dijumlah} (terlalu banyak pohon \emph{bisa} overfitting, lihat Diagram~\ref{dg:boostkurva}).
\end{jebak}

% ============================================================ BAB 3
\rb{Model Linear dan Model Jarak}
{\lgd{biru2}{kelas 0}\lgd{merah3}{kelas 1}\lgd{biru1}{bobot dan skor}}
{Cepat, mudah dijelaskan, dan baseline yang kuat. Andalan untuk teks dan embedding karena fiturnya sangat banyak.}

\subsection{Regresi logistik (dan Ridge, LinearSVC)}
\begin{inti}[biru1]
Tiap fitur diberi \textbf{bobot}, lalu dijumlah jadi satu \textbf{skor}. Untuk klasifikasi, skor diubah jadi peluang 0--1 lewat
kurva S (sigmoid). Batas keputusan = tempat skor $= 0$ (peluang 0,5), dan bentuknya \textbf{selalu garis lurus} (bidang datar
kalau fiturnya banyak).
\end{inti}

\rumus[skor]{z = b_1 x_1 + b_2 x_2 + \dots + c} \rumus[sigmoid]{p = \frac{1}{1 + e^{-z}}}

\begin{diag}[dg:logreg]{Garis batas dan kurva S}
\figs[1.0]{logreg}
\textbf{Cara membaca:} kiri, garis hitam = semua titik yang skornya 0. Hasil latihan di sini: $b_1 = \lrwa$, $b_2 = \lrwb$,
$c = \lrc$, akurasi \lracc\,\%. Kanan, tiap titik diletakkan menurut skornya; makin jauh dari 0, makin yakin. Titik biru
yang jatuh di skor positif (atau jingga di skor negatif) adalah yang salah tebak.
\end{diag}

\begin{kenapa}[bagaimana bobot dicari]
(1) Mulai dari bobot sembarang. (2) Hitung kesalahan (\emph{log-loss}: hukuman besar kalau yakin tapi salah). (3) Geser bobot
sedikit ke arah yang mengurangi kesalahan (\emph{gradient descent}). (4) Ulangi ratusan kali. Sklearn melakukannya otomatis saat
\kode{fit}; \kode{max\_iter=2000} = batas jumlah langkah.
\end{kenapa}

\begin{luar}
\kode{Ridge} dan \kode{Lasso} = versi regresi (tebak angka); Lasso bisa membuat bobot fitur tak berguna jadi tepat 0.
\kode{LinearSVC} dan \kode{SGDClassifier} = garis lurus juga, cara melatih berbeda. \kode{SVC} dengan kernel bisa melengkung karena
ruang datanya ``ditekuk'' dulu. Kenop bersama: \kode{C} (atau \kode{alpha} di Ridge). \kode{C} kecil = bobot ditahan, lebih aman dari
overfitting.
\end{luar}

\subsection{K-Nearest Neighbors (KNN)}
\begin{inti}[biru1]
Tidak ada pelatihan. Untuk menebak titik baru, cari \textbf{$k$ tetangga terdekat} di data latih, lalu ambil suara terbanyak
(atau rata-rata untuk regresi).
\end{inti}

\begin{diag}[dg:knn]{Pengaruh $k$}
\figs[1.0]{knn}
Bintang = titik yang ditanya; garis = tetangga yang ikut memilih; lingkaran putus-putus = batas tetangga terjauh.
\textbf{Cara membaca:} $k = 1$ (kiri) membuat batas berliku-liku dan pulau-pulau kecil karena satu titik aneh langsung berpengaruh.
$k = 15$ (kanan) membuat batas halus. Tebakan untuk bintang yang sama pun bisa berubah.
\end{diag}

\begin{jebak}
Model berbasis jarak (KNN, SVC, KMeans, juga regresi logistik saat dilatih) \textbf{wajib} diberi \kode{StandardScaler()} dulu.
Tanpa itu, kolom ``gaji'' (jutaan) menenggelamkan kolom ``umur'' (puluhan) karena jaraknya dihitung dari angka mentah.
Model pohon tidak butuh scaler.
\end{jebak}

% ============================================================ BAB 4
\rb{Neural Network dan Transfer Learning}
{\lgd{biru1}{input}\lgd{hijau1}{lapisan tersembunyi}\lgd{merah1}{output}}
{Dipakai untuk teks, gambar, audio, dan deret waktu. Di toolkit hampir selalu memakai model yang sudah dilatih orang lain.}

\begin{diag}[dg:nn]{MLP: regresi logistik yang disusun berlapis}
\figs[0.72]{nn}
Setiap lingkaran = satu neuron = jumlah berbobot dari lapisan sebelumnya, lalu ReLU (angka negatif dijadikan 0). Setiap garis = satu
bobot. \textbf{Cara membaca:} lapisan awal mendeteksi pola sederhana, lapisan berikutnya menggabungkannya jadi pola rumit. Saat
\kode{fit}, semua bobot diperbaiki bersama (\emph{backpropagation}), mirip gradient descent di bab regresi logistik.
\end{diag}

{\small
\begin{tabularx}{\linewidth}{@{}>{\bfseries\raggedright\arraybackslash}p{3.2cm}XP{4.2cm}@{}}
\toprule
Jenis & Idenya & Notebook\\ \midrule
\cl{hijau1}{MLP} & seperti diagram di atas; untuk tabel atau embedding. \kode{hidden\_layer\_sizes=(256,)} = satu lapisan 256 neuron
& \kode{tabular/deep\_klasifikasi}\\
\cl{hijau1}{CNN (ResNet)} & filter kecil 3$\times$3 digeser ke seluruh gambar mencari garis $\to$ sudut $\to$ tekstur $\to$ bentuk
& \kode{medis/klasifikasi\_citra}\\
\cl{hijau1}{YOLO} & CNN yang langsung menebak kotak + kelas objek dalam satu kali lihat & \kode{vision/deteksi}\\
\cl{hijau1}{LSTM} & membaca urutan satu per satu sambil membawa ``ingatan'' & \kode{timeseries/deep\_lstm}\\
\cl{hijau1}{Transformer} & tiap kata ``melihat'' semua kata lain sekaligus (\emph{attention}); IndoBERT, Chronos
& \kode{text/klasifikasi\_finetune}\\ \bottomrule
\end{tabularx}}

\begin{inti}[hijau1]
\textbf{Transfer learning.} Melatih neural network dari nol butuh jutaan data dan GPU berhari-hari. Toolkit memakai model
\emph{pretrained}, lalu salah satu dari dua cara: (1) \textbf{ambil embedding} (ringkasan berupa ratusan angka) dan latih
LogisticRegression di atasnya: cepat dan murah; (2) \textbf{fine-tune}: lanjutkan latihan sebentar dengan data lomba: lebih akurat,
butuh GPU.
\end{inti}

\begin{soal}{text/klasifikasi\_embedding (cara 1)}
\begin{lstlisting}
encoder = SentenceTransformer(EMBEDDING_MODEL)       # model pretrained, tidak dilatih ulang
X_train = encoder.encode(train_texts)                # tiap kalimat -> ratusan angka
model = LogisticRegression(C=1.0, max_iter=2000)     # model biasa di atas embedding
\end{lstlisting}
\end{soal}

\begin{kenapa}[kenapa embedding berguna]
Model pretrained sudah membaca jutaan kalimat atau gambar, sehingga ia ``tahu'' bahwa \emph{murah} dan \emph{terjangkau}
mirip, atau bahwa telinga runcing + kumis = kucing. Embedding menyimpan pengetahuan itu sebagai posisi: yang mirip
berdekatan. Model sederhana di atasnya tinggal menarik garis.
\end{kenapa}

% ============================================================ BAB 5
\rb{Tanpa Label: Clustering, Reduksi Dimensi, Anomali}
{\lgd{biru2}{data}\lgd{merah3}{hasil proyeksi / titik terpilih}\lgd{biru3}{potongan acak}}
{Tidak ada kolom $y$. Model mencari struktur sendiri; kita yang menafsirkan hasilnya.}

\subsection{KMeans}
\begin{inti}[biru1]
Taruh $k$ pusat secara acak, lalu ulangi dua langkah: \textbf{(1)} tiap titik ikut pusat terdekat, \textbf{(2)} tiap pusat pindah
ke tengah (rata-rata) anggotanya. Berhenti saat tidak ada titik yang pindah kelompok.
\end{inti}

\begin{diag}[dg:kmeans]{KMeans langkah demi langkah ($k = 4$)}
\figs[0.82]{kmeans}
$\times$ = pusat klaster. \textbf{Cara membaca:} panel 0, pusat diambil dari titik acak (dua pusat bisa jatuh di kelompok yang
sama). Panel 1, garis tipis menunjukkan tiap titik memilih pusat terdekat. Panel 2, pusat bergeser ke tengah anggotanya.
Setelah \kmiter{} putaran, hasilnya stabil.
\end{diag}

\rumus[inertia]{\text{inertia} = \sum_{\text{titik}} \text{jarak}(\text{titik},\ \text{pusatnya})^2 \quad \text{(makin kecil = makin rapat)}}

\begin{jebak}
KMeans \textbf{selalu} memberi jawaban, bahkan untuk $k$ yang salah. Memilih $k$ adalah tugas kita: grafik \emph{elbow}
(inertia vs $k$, cari ``siku'') atau \emph{silhouette score} (makin tinggi makin baik). Skalakan fitur dulu. Nama klaster
(``pembeli hemat'') kita yang memberi, dengan melihat rata-rata fitur tiap klaster.
\end{jebak}

\begin{luar}
\kode{DBSCAN}/\kode{HDBSCAN}: klaster = daerah padat, jumlahnya otomatis, titik sepi diberi label noise ($-1$); cocok untuk bentuk
tak bulat (GPS, trajektori). \kode{AgglomerativeClustering}: gabung dua kelompok terdekat berulang-ulang (dendrogram).
\kode{GaussianMixture}: seperti KMeans tetapi tiap titik punya peluang ikut tiap klaster.
\end{luar}

\subsection{PCA, UMAP, t-SNE}
\begin{inti}[biru1]
Meringkas banyak kolom jadi sedikit kolom baru. PCA mencari \textbf{arah tempat data paling menyebar}, lalu memproyeksikan
(menjatuhkan) setiap titik ke arah itu. Komponen kedua = arah tersebar berikutnya yang tegak lurus dengan yang pertama.
\end{inti}

\begin{diag}[dg:pca]{Arah yang baik vs arah yang buruk}
\figs[1.0]{pca}
Garis tipis = jalur jatuhnya titik ke garis biru; titik jingga = posisi baru. Strip di bawah = data setelah tinggal 1 kolom.
\textbf{Cara membaca:} di arah PCA (kiri) titik-titik tetap terpisah lebar, \pcasatu\,\% sebaran tersimpan. Di arah tegak lurus
(kanan) semuanya menumpuk, hanya \pcadua\,\% tersimpan. Di sklearn angka ini ada di \kode{explained\_variance\_ratio\_}.
\end{diag}

\begin{luar}
\textbf{UMAP} dan \textbf{t-SNE} hanya untuk \emph{melihat}. Keduanya menjaga agar tetangga dekat tetap dekat, jadi kelompok tampak
jelas di gambar 2D, tetapi jarak antar-kelompok dan ukuran kelompok di gambar tidak bermakna. PCA juga dipakai untuk
mengecilkan embedding (misal 768 $\to$ 50 kolom) sebelum clustering.
\end{luar}

\subsection{Isolation Forest}
\begin{inti}[merah1]
Potong data secara acak berulang-ulang. Titik aneh yang jauh dari kerumunan \textbf{cepat terpisah sendirian}; titik normal di
tengah kerumunan butuh banyak potongan. \textbf{Sedikit potongan = anomali.}
\end{inti}

\begin{diag}[dg:iforest]{Berapa potongan untuk mengisolasi satu titik}
\figs[1.0]{iforest}
Cincin jingga = titik yang sedang diisolasi; garis = potongan acak; kotak berarsir = kotak terakhir yang hanya berisi titik itu.
\textbf{Cara membaca:} rata-rata 300 kali percobaan, titik terpencil butuh \ifluar{} potongan, titik di tengah butuh \iftengah{}
potongan. IsolationForest mengulang ini dengan ratusan pohon acak dan mengubah rata-rata panjang jalur jadi skor anomali.
\end{diag}

\begin{luar}
\kode{LocalOutlierFactor}: titik yang daerahnya jauh lebih sepi daripada daerah tetangganya. \kode{ECOD}, \kode{HBOS} (pyod): seberapa
ekstrem nilai tiap kolom. Kenop \kode{contamination} = perkiraan persen data yang anomali; ini tebakan kita, jadi selalu cek
baris yang ditandai secara manual (\kode{tabular/cek\_outlier\_input} menyimpannya ke \kode{baris\_curiga.csv}).
\end{luar}

% ============================================================ BAB 6
\rb{Teks: Dari Kata Menjadi Angka}
{\lgd{biru1}{nilai TF-IDF}}
{Model hanya mengerti angka. Langkah pertama setiap notebook teks adalah mengubah kalimat menjadi tabel angka.}

\begin{inti}[biru1]
\textbf{TF-IDF} memberi nilai tinggi pada kata yang \textbf{sering di dokumen ini} (TF) tetapi \textbf{jarang di dokumen lain}
(IDF). Kata yang muncul di mana-mana (``dan'', ``barang'') bernilai kecil karena tidak membedakan apa-apa. Setiap kata menjadi
satu kolom; LogisticRegression/LinearSVC lalu memberi bobot ke tiap kolom.
\end{inti}

\rumus[seperti sklearn, smooth\_idf]{\text{TF} = \frac{\text{jumlah kata di dokumen}}{\text{panjang dokumen}}, \qquad
\text{IDF} = \ln\frac{1 + N}{1 + \text{df}} + 1, \qquad \text{TF-IDF} = \text{TF} \times \text{IDF}}

\begin{soal}{4 ulasan, nilai kata di dokumen 4 ``barang bagus harga murah murah''}
Dokumen: (1) pengiriman cepat dan barang bagus; (2) barang rusak dan pengiriman lambat sekali; (3) pelayanan ramah dan barang
bagus sekali; (4) barang bagus harga murah murah. $N = 4$; df = berapa dokumen memuat kata itu.\par\smallskip
\begin{center}
\begin{tabular}{@{}lccccl@{}}
\toprule
kata & TF & df & IDF & TF-IDF & \\ \midrule
\textbf{murah} & 2/5 = 0{,}40 & 1 dari 4 & 1{,}92 & 0{,}767 & \tikz\fill[biru1] (0,0) rectangle (2.60cm,0.18cm);\\
\textbf{harga} & 1/5 = 0{,}20 & 1 dari 4 & 1{,}92 & 0{,}383 & \tikz\fill[biru1] (0,0) rectangle (1.30cm,0.18cm);\\
\textbf{bagus} & 1/5 = 0{,}20 & 3 dari 4 & 1{,}22 & 0{,}245 & \tikz\fill[biru1] (0,0) rectangle (0.83cm,0.18cm);\\
\textbf{barang} & 1/5 = 0{,}20 & 4 dari 4 & 1{,}00 & 0{,}200 & \tikz\fill[biru1] (0,0) rectangle (0.68cm,0.18cm);\\
\bottomrule

\end{tabular}
\end{center}
``murah'' paling tinggi: muncul dua kali dan hanya di dokumen ini. ``barang'' paling rendah: ada di semua dokumen.
\end{soal}

\begin{luar}
\kode{ngram\_range=(1,2)}: pasangan kata ikut dihitung, sehingga ``tidak bagus'' jadi satu kolom. \kode{min\_df=2}: buang kata yang
muncul di kurang dari 2 dokumen. \kode{analyzer="char\_wb"}: potongan huruf, tahan salah ketik dan bahasa gaul.
\kode{sublinear\_tf=True}: kata yang muncul 10 kali tidak dianggap 10 kali lebih penting.
\end{luar}

\sub{Embedding dan cosine similarity}
TF-IDF tidak tahu bahwa ``murah'' dan ``terjangkau'' mirip. Embedding (SentenceTransformer) tahu, karena model pretrained-nya sudah
membaca jutaan kalimat. Kemiripan dua kalimat diukur dari sudut antara dua vektornya:
\rumus[cosine similarity]{\cos\theta = \frac{\vec a \cdot \vec b}{\lVert \vec a\rVert\, \lVert \vec b\rVert} \qquad
(1 = \text{searah / sangat mirip},\ 0 = \text{tidak berhubungan})}
Dipakai di \kode{text/retrieval} (cari dokumen paling mirip pertanyaan), \kode{text/rag\_llm}, dan \kode{preprocessing/peta\_semantik}.

\begin{jebak}
Jangan \kode{fit} \kode{TfidfVectorizer} pada gabungan latih + test lalu melaporkan skor validasi: IDF ikut ``melihat'' data
validasi. Masukkan vectorizer ke \kode{Pipeline} atau \kode{fit} pada data latih saja, lalu \kode{transform} yang lain.
\end{jebak}

% ============================================================ BAB 7
\rb{Deret Waktu: Meramal dari Masa Lalu}
{\lgd{biru1}{$\alpha$ kecil}\lgd{merah3}{$\alpha$ besar}\lgd{merah1}{kejadian khusus}}
{Masa depan ditebak dari masa lalu. Tiga cara di toolkit: smoothing/SARIMAX, fitur lag + model tabel, dan model besar (Chronos).}

\subsection{Exponential smoothing}
\begin{inti}[biru1]
Rata-rata yang lebih berat ke data terbaru. \textbf{Level} = perkiraan ``nilai normal saat ini''; ramalan ke depan = level terakhir.
$\alpha$ mengatur seberapa cepat level mengikuti data baru.
\end{inti}

\rumus[level]{\ell_t = \alpha\, y_t + (1 - \alpha)\, \ell_{t-1}, \qquad \hat y_{t+1} = \ell_t}

\begin{diag}[dg:es]{Pengaruh $\alpha$}
\figs[0.9]{es}
\textbf{Cara membaca:} $\alpha = 0{,}1$ (biru) tenang tetapi selalu tertinggal di belakang data yang terus naik, MAE \esmaekecil.
$\alpha = 0{,}8$ (jingga) cepat mengikuti, tetapi ikut melonjak di bulan 22, MAE \esmaebesar. Ramalan keduanya datar karena model
ini tidak mengenal tren.
\end{diag}

\begin{kenapa}[kenapa ada Holt dan Holt-Winters]
Simple smoothing hanya punya level, jadi ramalannya datar. \textbf{Holt} menambah komponen \emph{tren} (naik berapa per bulan),
\textbf{Holt-Winters} menambah komponen \emph{musim} (pola yang berulang tiap 12 bulan atau tiap 7 hari). Di statsmodels:
\kode{ExponentialSmoothing(trend="add", seasonal="add", seasonal\_periods=12)}.
\end{kenapa}

\subsection{SARIMAX}
{\small
\begin{tabularx}{\linewidth}{@{}>{\bfseries}p{1.6cm}X@{}}
\toprule
AR & \emph{autoregressive}: nilai sekarang mirip nilai beberapa periode lalu\\
I & \emph{integrated}: modelkan selisih ($y_t - y_{t-1}$), bukan nilai mentah, supaya tren hilang\\
MA & \emph{moving average}: koreksi dari kesalahan ramalan periode lalu\\
S & \emph{seasonal}: AR/I/MA yang sama tetapi untuk jarak satu musim (12 bulan lalu)\\
X & \emph{exogenous}: faktor luar yang diketahui (libur, promo, harga)\\ \bottomrule
\end{tabularx}}

\subsection{Fitur lag + model tabel}
\begin{inti}[hijau1]
Cara yang sering paling kuat di lomba: ubah deret jadi \textbf{tabel biasa}. Tiap baris berisi nilai bulan lalu (\kode{lag\_1}),
tahun lalu (\kode{lag\_12}), rata-rata beberapa bulan sebelumnya, bulan ke berapa, lalu latih RandomForest/LightGBM biasa.
\end{inti}
\begin{center}
\small
\begin{tabular}{@{}lcccc@{}}
\toprule
bulan & jual ($y$) & \kode{lag\_1} & \kode{lag\_12} & \kode{rata\_3}\\ \midrule
2024-01 & 120 & -- & -- & --\\
2024-02 & 132 & 120 & -- & --\\
2024-03 & 128 & 132 & -- & --\\
2024-04 & 140 & 128 & -- & 126{,}7\\
\dots & \dots & \dots & \dots & \dots\\
2025-01 & 151 & 145 & 120 & 143{,}0\\ \bottomrule
\end{tabular}
\end{center}

\begin{jebak}
(1) Deret waktu \textbf{tidak boleh diacak} saat split; validasi harus sesudah data latih (\kode{TimeSeriesSplit},
Diagram~\ref{dg:kfold}). (2) \kode{rata\_3} harus dihitung dari bulan \emph{sebelumnya} (\kode{shift(1)} dulu, baru
\kode{rolling(3)}); kalau bulan ini ikut dihitung, model mengintip jawaban. (3) Saat meramal beberapa langkah, \kode{lag\_1} bulan
depan belum diketahui; pakai hasil ramalan sebelumnya atau model terpisah per horizon.
\end{jebak}

% ============================================================ BAB 8
\rb{Rekomendasi dan Graf}
{\lgd{biru1}{tebakan model}\lgd{biru2}{komunitas 1}\lgd{merah3}{komunitas 2}\lgd{hijau2}{komunitas 3}}
{Dua jenis data yang bentuknya bukan tabel biasa: matriks orang $\times$ barang, dan jaringan simpul + sisi.}

\subsection{Matrix factorization (SVD)}
\begin{inti}[biru1]
Tiap orang dan tiap film diringkas jadi beberapa angka ``selera'' (faktor). Rating tebakan = rata-rata umum + kebiasaan orang
(pelit/murah memberi nilai) + kebiasaan film (memang disukai/tidak) + kecocokan selera orang dan film.
\end{inti}

\rumus[seperti \kode{surprise.SVD}]{\hat r_{u,i} = \mu + b_u + b_i + \vec p_u \cdot \vec q_i}

\begin{soal}{5 orang $\times$ 6 film, 2 faktor (dilatih dengan SGD di \kode{buat\_gambar.py})}
\begin{center}
\begin{tabular}{@{}lcccccc@{}}
\toprule
 & Horor A & Horor B & Komedi A & Komedi B & Drama A & Drama B\\ \midrule
\textbf{Ani} & 5 & 4 & 1 & \textcolor{biru1}{\textbf{\emph{1{,}4}}} & 2 & \textcolor{biru1}{\textbf{\emph{1{,}6}}}\\
\textbf{Budi} & 4 & \textcolor{biru1}{\textbf{\emph{3{,}9}}} & 2 & 1 & \textcolor{biru1}{\textbf{\emph{2{,}4}}} & 2\\
\textbf{Citra} & 1 & 2 & 5 & 4 & \textcolor{biru1}{\textbf{\emph{4{,}3}}} & 4\\
\textbf{Dodi} & \textcolor{biru1}{\textbf{\emph{1{,}4}}} & 1 & 4 & 5 & 4 & \textcolor{biru1}{\textbf{\emph{3{,}5}}}\\
\textbf{Eka} & 2 & \textcolor{biru1}{\textbf{\emph{3{,}2}}} & \textcolor{biru1}{\textbf{\emph{5{,}0}}} & 4 & 5 & 5\\
\bottomrule

\end{tabular}
\end{center}
Angka tegak = rating asli; angka \textcolor{biru1}{\textbf{\emph{miring biru}}} = tebakan untuk sel kosong. Contoh Ani $\times$
Komedi B: $\mu + b_u + b_i + \vec p\cdot\vec q = \mfmu + (\mfbu) + (\mfbi) + (\mfpq) = \mfhasil$. Ani suka horor dan memberi
Komedi A nilai 1, jadi kecocokan seleranya negatif.
\end{soal}

\begin{luar}
\kode{KNNBaseline} (orang/film yang mirip), \kode{SVDpp}, \kode{NMF}. Untuk data klik tanpa rating: \kode{ALS} dan \kode{BPR}
(\kode{recsys/top\_n}). Baseline \textbf{popularitas} (rekomendasikan yang paling banyak dibeli) sering sulit dikalahkan; di data demo
toolkit ia justru menang (\kode{KATALOG.md}).
\end{luar}

\subsection{Deteksi komunitas (Louvain)}
\begin{inti}[biru1]
Cari kelompok simpul yang \textbf{banyak saling terhubung di dalam} dan \textbf{sedikit ke luar}. Louvain memindah simpul antar
kelompok selama skor \emph{modularitas} naik, lalu menggabungkan tiap kelompok jadi satu simpul besar dan mengulang.
\end{inti}
\begin{center}
\begin{minipage}{0.42\linewidth}\centering\begin{tikzpicture}[x=3.2cm,y=3.2cm]
\draw[black!30,line width=0.6pt] (0.7500,0.6200)--(0.6500,0.7499);
\draw[black!30,line width=0.6pt] (0.7500,0.6200)--(0.4500,0.7499);
\draw[black!30,line width=0.6pt] (0.7500,0.6200)--(0.3500,0.6200);
\draw[black!30,line width=0.6pt] (0.7500,0.6200)--(0.4500,0.4901);
\draw[black!30,line width=0.6pt] (0.6500,0.7499)--(0.4500,0.7499);
\draw[black!30,line width=0.6pt] (0.6500,0.7499)--(0.3500,0.6200);
\draw[black!30,line width=0.6pt] (0.6500,0.7499)--(0.4500,0.4901);
\draw[black!30,line width=0.6pt] (0.6500,0.7499)--(0.6500,0.4901);
\draw[black!30,line width=0.6pt] (0.4500,0.7499)--(0.3500,0.6200);
\draw[black!30,line width=0.6pt] (0.4500,0.7499)--(0.4500,0.4901);
\draw[black!30,line width=0.6pt] (0.3500,0.6200)--(0.4500,0.4901);
\draw[black!30,line width=0.6pt] (1.5581,0.7862)--(1.2502,0.6671);
\draw[black!30,line width=0.6pt] (1.5581,0.7862)--(1.3419,0.5338);
\draw[black!30,line width=0.6pt] (1.5581,0.7862)--(1.6498,0.6529);
\draw[black!30,line width=0.6pt] (1.3583,0.7933)--(1.2502,0.6671);
\draw[black!30,line width=0.6pt] (1.3583,0.7933)--(1.3419,0.5338);
\draw[black!30,line width=0.6pt] (1.3583,0.7933)--(1.5417,0.5267);
\draw[black!30,line width=0.6pt] (1.3583,0.7933)--(1.6498,0.6529);
\draw[black!30,line width=0.6pt] (1.2502,0.6671)--(1.5417,0.5267);
\draw[black!30,line width=0.6pt] (1.2502,0.6671)--(1.6498,0.6529);
\draw[black!30,line width=0.6pt] (1.3419,0.5338)--(1.6498,0.6529);
\draw[black!30,line width=0.6pt] (1.5417,0.5267)--(1.6498,0.6529);
\draw[black!30,line width=0.6pt] (0.9168,0.3364)--(0.8009,0.2141);
\draw[black!30,line width=0.6pt] (0.9168,0.3364)--(0.8841,0.0777);
\draw[black!30,line width=0.6pt] (0.9168,0.3364)--(1.0832,0.0636);
\draw[black!30,line width=0.6pt] (0.9168,0.3364)--(1.1159,0.3223);
\draw[black!30,line width=0.6pt] (0.8009,0.2141)--(0.8841,0.0777);
\draw[black!30,line width=0.6pt] (0.8009,0.2141)--(1.1991,0.1859);
\draw[black!30,line width=0.6pt] (0.8841,0.0777)--(1.0832,0.0636);
\draw[black!30,line width=0.6pt] (0.8841,0.0777)--(1.1991,0.1859);
\draw[black!30,line width=0.6pt] (1.0832,0.0636)--(1.1991,0.1859);
\draw[black!30,line width=0.6pt] (0.7500,0.6200)--(1.3583,0.7933);
\draw[black!30,line width=0.6pt] (0.3500,0.6200)--(0.8009,0.2141);
\draw[black!30,line width=0.6pt] (1.3419,0.5338)--(1.0832,0.0636);
\draw[black!30,line width=0.6pt] (0.6500,0.4901)--(1.1991,0.1859);
\fill[biru2,draw=white,line width=0.6pt] (0.7500,0.6200) circle (3.2pt);
\fill[biru2,draw=white,line width=0.6pt] (0.6500,0.7499) circle (3.2pt);
\fill[biru2,draw=white,line width=0.6pt] (0.4500,0.7499) circle (3.2pt);
\fill[biru2,draw=white,line width=0.6pt] (0.3500,0.6200) circle (3.2pt);
\fill[biru2,draw=white,line width=0.6pt] (0.4500,0.4901) circle (3.2pt);
\fill[biru2,draw=white,line width=0.6pt] (0.6500,0.4901) circle (3.2pt);
\fill[merah3,draw=white,line width=0.6pt] (1.5581,0.7862) circle (3.2pt);
\fill[merah3,draw=white,line width=0.6pt] (1.3583,0.7933) circle (3.2pt);
\fill[merah3,draw=white,line width=0.6pt] (1.2502,0.6671) circle (3.2pt);
\fill[merah3,draw=white,line width=0.6pt] (1.3419,0.5338) circle (3.2pt);
\fill[merah3,draw=white,line width=0.6pt] (1.5417,0.5267) circle (3.2pt);
\fill[merah3,draw=white,line width=0.6pt] (1.6498,0.6529) circle (3.2pt);
\fill[hijau2,draw=white,line width=0.6pt] (0.9168,0.3364) circle (3.2pt);
\fill[hijau2,draw=white,line width=0.6pt] (0.8009,0.2141) circle (3.2pt);
\fill[hijau2,draw=white,line width=0.6pt] (0.8841,0.0777) circle (3.2pt);
\fill[hijau2,draw=white,line width=0.6pt] (1.0832,0.0636) circle (3.2pt);
\fill[hijau2,draw=white,line width=0.6pt] (1.1991,0.1859) circle (3.2pt);
\fill[hijau2,draw=white,line width=0.6pt] (1.1159,0.3223) circle (3.2pt);
\end{tikzpicture}
\end{minipage}\hfill
\begin{minipage}{0.54\linewidth}\small
\rumus[modularitas]{Q = \sum_c \left[\frac{L_c}{m} - \left(\frac{d_c}{2m}\right)^2\right]}
$L_c$ = sisi di dalam kelompok $c$; $d_c$ = jumlah derajat simpul di $c$; $m$ = total sisi (\jumlahsisi).
Pembagian di gambar: $Q = \qgrup$. Semua dijadikan satu kelompok: $Q = \qsatu$. Makin tinggi (maks.\ 1) = kelompok makin jelas.
\end{minipage}
\end{center}

\begin{luar}
Fitur graf untuk model tabel: \textbf{degree} (jumlah teman), \textbf{PageRank} (dihubungi simpul penting), \textbf{betweenness}
(sering jadi jembatan). Fitur ini lalu masuk LightGBM biasa (\kode{graph/klasifikasi\_node}). \textbf{GCN}: neural network yang
membuat tiap simpul mengumpulkan informasi dari tetangganya, beberapa lapis (\kode{graph/gcn}).
\end{luar}

% ============================================================ BAB 9
\rb{Menilai dan Memeras Skor}
{\lgd{biru2}{latih / kelas 0}\lgd{merah3}{validasi / kelas 1}\lgd{hijau1}{tebakan benar}\lgd{merah1}{tebakan salah}}
{Skor validasi menentukan semua keputusan. Kalau cara menghitungnya salah, semua keputusan ikut salah.}

\subsection{Validasi silang (KFold)}
\begin{inti}[biru1]
Satu kali split bisa beruntung atau sial. KFold membagi data jadi 5 bagian; model dilatih 5 kali, tiap kali satu bagian berbeda
jadi validasi. Skor akhir = rata-rata 5 skor, dan simpangannya menunjukkan seberapa stabil model.
\end{inti}

\begin{diag}[dg:kfold]{KFold vs TimeSeriesSplit}
\figs[0.85]{kfold}
\textbf{Cara membaca:} kiri, setiap bagian kebagian jadi validasi tepat satu kali. Kanan, validasi selalu \emph{sesudah} data
latih dan data latih makin panjang; bagian sesudah validasi tidak dipakai supaya model tidak mengintip masa depan.
\end{diag}

{\small
\begin{tabularx}{\linewidth}{@{}>{\bfseries}p{3.2cm}X@{}}
\toprule
StratifiedKFold & proporsi kelas di tiap lipatan dijaga sama; wajib bila kelas tidak seimbang (95\,\% normal, 5\,\% fraud)\\
GroupKFold & satu orang/pasien/toko tidak boleh muncul di latih dan validasi sekaligus\\
TimeSeriesSplit & untuk deret waktu; validasi selalu sesudah latih\\ \bottomrule
\end{tabularx}}

\subsection{Precision, recall, F1, dan threshold}
\begin{inti}[biru1]
Classifier mengeluarkan \textbf{peluang}. Kita memilih ambang (\emph{threshold}): di atasnya dianggap kelas 1. Menggeser ambang
\textbf{menukar} dua jenis kesalahan: alarm palsu (FP) dan yang terlewat (FN).
\end{inti}

\begin{diag}[dg:threshold]{Satu set prediksi, threshold 0,5}
\figs[0.9]{threshold}
Titik penuh = tebakan benar; cincin kosong = tebakan salah. \textbf{Cara membaca:} cincin biru di daerah kanan = FP (aslinya 0,
ditebak 1); cincin jingga di kiri = FN (aslinya 1, ditebak 0). Geser garis ke kiri: FN berkurang, FP bertambah.
\end{diag}

\rumus{\text{precision} = \frac{TP}{TP + FP}, \qquad \text{recall} = \frac{TP}{TP + FN}, \qquad
F_1 = \frac{2 \cdot \text{precision} \cdot \text{recall}}{\text{precision} + \text{recall}}}

\begin{center}
\small
\begin{tabular}{@{}cccccccc@{}}
\toprule
threshold & TP & FP & FN & TN & precision & recall & F1\\ \midrule
0{,}3 & 25 & 15 & 0 & 30 & 0{,}62 & 1{,}00 & \textbf{0{,}77}\\
0{,}5 & 22 & 6 & 3 & 39 & 0{,}79 & 0{,}88 & \textbf{0{,}83}\\
0{,}7 & 14 & 0 & 11 & 45 & 1{,}00 & 0{,}56 & \textbf{0{,}72}\\
\bottomrule

\end{tabular}
\end{center}

{\small
\begin{tabularx}{\linewidth}{@{}>{\bfseries}p{2.4cm}X@{}}
\toprule
F1 macro & F1 dihitung per kelas lalu dirata-rata; kelas kecil sama penting dengan kelas besar\\
AUC & kualitas \emph{urutan} peluang (apakah kelas 1 umumnya diberi peluang lebih tinggi); tidak bergantung threshold\\
MAE & rata-rata $|y - \hat y|$; mudah dibaca dalam satuan aslinya\\
RMSE & $\sqrt{\text{rata-rata }(y - \hat y)^2}$; kesalahan besar dihukum lebih berat\\
$R^2$ & 1 = sempurna, 0 = setara menebak rata-rata, negatif = lebih buruk dari rata-rata\\ \bottomrule
\end{tabularx}}

\begin{jebak}
Pakai metrik yang \textbf{persis sama} dengan yang ditulis di soal. Akurasi 95\,\% bisa tidak berarti apa-apa kalau 95\,\% data
memang kelas 0 (model yang selalu menjawab 0 pun dapat 95\,\%). Untuk F1, threshold 0,5 belum tentu terbaik; cari di data
validasi (\kode{tabular/imbalanced}), jangan di data test.
\end{jebak}

\subsection{Ensembling, tuning, dan SHAP}
{\small
\begin{tabularx}{\linewidth}{@{}>{\bfseries\raggedright\arraybackslash}p{2.8cm}>{\raggedright\arraybackslash}XP{4.6cm}@{}}
\toprule
Ensembling / blending & rata-ratakan peluang beberapa model berbeda jenis (misal LightGBM 0,62, CatBoost 0,48, RandomForest
0,71 $\to$ 0,60). Kesalahan mereka tidak sama, jadi saling menutupi & \kode{tabular/ensembling\_klasifikasi}, \kode{tabular/blending}\\
Tuning (Optuna) & mencoba kombinasi kenop dan belajar dari percobaan sebelumnya: area yang skornya bagus dicoba lebih sering.
Lebih hemat daripada mencoba semua kombinasi & \kode{tabular/tuning\_klasifikasi}\\ \bottomrule
\end{tabularx}}

\begin{diag}[dg:shap]{SHAP: kenapa pelanggan ini ditebak 59\,\% akan berhenti}
\figs[0.8]{shap}
\textbf{Cara membaca:} mulai dari rata-rata tebakan model untuk semua pelanggan (0,30). Tiap fitur pelanggan ini mendorong
naik (\textcolor{merah3}{jingga}) atau turun (\textcolor{biru2}{biru}); jumlah semua dorongan = prediksi akhir 0,59. Notebook:
\kode{tabular/explain\_klasifikasi}, \kode{visualisasi/penjelasan\_model}.
\end{diag}

\begin{kenapa}[ensembling vs tuning, mana dulu]
Tuning biasanya menaikkan skor sedikit (sering $<$ 1 poin) dan makan waktu. Ensembling 2--3 model yang sudah lumayan dan
berbeda jenis sering memberi kenaikan serupa dengan usaha lebih kecil. Urutan hemat waktu saat lomba: baseline $\to$ fitur
yang lebih baik $\to$ ganti model $\to$ ensembling $\to$ tuning.
\end{kenapa}

% ============================================================ LEMBAR RINGKAS
\clearpage
\section*{Lembar ringkas satu halaman}
\addcontentsline{toc}{section}{Lembar ringkas satu halaman}
{\small
\begin{tabularx}{\linewidth}{@{}>{\bfseries\raggedright\arraybackslash}p{3.3cm}XP{4.0cm}@{}}
\toprule
Metode & Idenya dalam satu kalimat & Kenop utama\\ \midrule
Decision Tree & tanya ya/tidak berulang, potong ruang dengan garis lurus & \kode{max\_depth}, \kode{min\_samples\_leaf}\\
Random Forest & ratusan pohon dalam dari sampel acak, dirata-rata & \kode{n\_estimators}, \kode{max\_features}\\
XGBoost/LightGBM & pohon dangkal berurutan, tiap pohon menebak sisa kesalahan & \kode{learning\_rate}, \kode{n\_estimators}, \kode{num\_leaves}\\
CatBoost & boosting yang jago kolom kategori & \kode{iterations}, \kode{depth}\\
LogisticRegression & jumlah berbobot $\to$ sigmoid $\to$ peluang; batas lurus & \kode{C}\\
Ridge / Lasso & regresi linear dengan rem bobot & \kode{alpha}\\
SVC & batas dengan jarak (margin) terbesar; kernel = bisa melengkung & \kode{C}, \kode{kernel}\\
KNN & suara $k$ tetangga terdekat & \kode{n\_neighbors}\\
MLP / CNN / Transformer & lapisan neuron; pretrained + embedding atau fine-tune & epoch, learning rate, batch size\\
KMeans & titik ikut pusat terdekat, pusat pindah ke tengah & \kode{n\_clusters}\\
DBSCAN & klaster = daerah padat; sisanya noise & \kode{eps}, \kode{min\_samples}\\
PCA & proyeksi ke arah sebaran terbesar & \kode{n\_components}\\
IsolationForest & sedikit potongan acak untuk terisolasi = anomali & \kode{contamination}\\
TF-IDF & sering di dokumen ini, jarang di lainnya & \kode{ngram\_range}, \kode{min\_df}\\
Exp.\ smoothing & rata-rata berat ke data terbaru & \kode{smoothing\_level} ($\alpha$)\\
SARIMAX & AR + selisih + koreksi error + musim + faktor luar & \kode{order}, \kode{seasonal\_order}\\
Fitur lag + LightGBM & deret waktu diubah jadi tabel & lag yang dipakai\\
SVD (recsys) & orang dan barang jadi vektor selera & \kode{n\_factors}, \kode{reg\_all}\\
Louvain & kelompok yang padat di dalam, jarang ke luar & \kode{resolution}\\ \bottomrule
\end{tabularx}}

\vspace{4pt}
{\small
\begin{tabularx}{\linewidth}{@{}P{4.2cm}X@{}}
\toprule
Rumus & \\ \midrule
Gini & $1 - \sum_k p_k^2$\\
Boosting & $F_m(x) = F_{m-1}(x) + \eta\, h_m(x)$, $h_m$ menebak $y - F_{m-1}$\\
Regresi logistik & $z = b_1x_1 + \dots + c$; \ $p = 1/(1+e^{-z})$\\
TF-IDF & $\text{TF} \times \left(\ln\frac{1+N}{1+\text{df}} + 1\right)$\\
Cosine similarity & $\vec a\cdot\vec b / (\lVert\vec a\rVert\lVert\vec b\rVert)$\\
Exponential smoothing & $\ell_t = \alpha y_t + (1-\alpha)\ell_{t-1}$\\
Matrix factorization & $\hat r = \mu + b_u + b_i + \vec p_u\cdot\vec q_i$\\
Modularitas & $\sum_c [L_c/m - (d_c/2m)^2]$\\
Precision / recall / F1 & $TP/(TP+FP)$; \ $TP/(TP+FN)$; \ $2PR/(P+R)$\\
RMSE / MAE & $\sqrt{\overline{(y-\hat y)^2}}$; \ $\overline{|y-\hat y|}$\\ \bottomrule
\end{tabularx}}

\section*{Kesalahan pemula yang perlu diingat}
\addcontentsline{toc}{section}{Kesalahan pemula yang perlu diingat}
{\small
\begin{tabularx}{\linewidth}{@{}P{4.4cm}>{\raggedright\arraybackslash}X>{\raggedright\arraybackslash}X@{}}
\toprule
Gejala & Penyebab & Perbaikan\\ \midrule
skor validasi 0,99, skor lomba jauh lebih rendah & kebocoran data (scaler/TF-IDF di-\kode{fit} pada semua data, kolom turunan target)
& \kode{Pipeline}; cek kolom satu per satu\\
KNN/SVC/KMeans hasilnya aneh & fitur tidak diskalakan & \kode{make\_pipeline(StandardScaler(), model)}\\
skor latih jauh di atas skor validasi & overfitting & kurangi \kode{max\_depth}/\kode{num\_leaves}, naikkan regularisasi\\
deret waktu skornya terlalu bagus & split diacak / rata-rata bergulir memuat bulan ini & \kode{TimeSeriesSplit}; \kode{shift(1)} sebelum \kode{rolling}\\
akurasi tinggi tetapi kelas kecil tidak pernah tertebak & kelas tidak seimbang & pakai F1 macro, \kode{class\_weight="balanced"}, atur threshold\\
\kode{learning\_rate} diturunkan, skor malah turun & jumlah pohon tidak dinaikkan & naikkan \kode{n\_estimators} bersamaan\\
skor berubah-ubah tiap \emph{Run All} & seed tidak dipasang & \kode{random\_state=SEED} di model dan split\\ \bottomrule
\end{tabularx}}

\end{document}
